% ---------------------------------------------------------------------------
% Chapter: Troubleshooting
% ---------------------------------------------------------------------------
\chapter{Troubleshooting}\label{chap:troubleshooting}
% Long monospaced identifiers leave little room for line breaking; allow some
% extra inter-word stretch in this chapter (the group keeps it local).
\begingroup\setlength{\emergencystretch}{3em}% chapter-local

This chapter collects the problems that users most often run into, grouped by
symptom, together with their causes and solutions. If a problem is not
covered here, run the failing code with warnings enabled
(\code{python -W default ...}), check the installation with the commands of
Section~\ref{sec:install-verify}, and report the issue with the full error
message and the output of \code{pip list} at
\url{https://github.com/EricDMW/My_Tool_Box/issues}.

\section{Rendering on machines without a display}\label{sec:trouble-headless}

\paragraph{Symptom.} Rendering fails or hangs on a compute cluster, in a
container, over SSH or in continuous integration, often with messages such as
\code{cannot connect to X server}, \code{no display name} or
\code{TclError}.

\paragraph{Solution.} Use \code{render\_mode="rgb\_array"}, which draws
off-screen and never needs a display: the matplotlib renderers draw on an Agg
canvas that is independent of the active matplotlib backend, and the
Pistonball renderer draws on pygame surfaces without initialising the display.
In addition, select non-interactive back ends for any other plotting code:

\begin{shellcode}
export MPLBACKEND=Agg            # matplotlib: off-screen rendering
export SDL_VIDEODRIVER=dummy     # pygame: no window system
python train.py
\end{shellcode}

The same settings can be made at the top of a script, before matplotlib or
pygame is imported:

\begin{pycode}
import os

os.environ.setdefault("MPLBACKEND", "Agg")
os.environ.setdefault("SDL_VIDEODRIVER", "dummy")
\end{pycode}

Frames returned by \code{env.render()} can then be recorded to GIF or MP4
(Section~\ref{sec:examples-recording}) and inspected later. The test suite,
the benchmark script and the CI workflow use exactly these settings.

\section{No window appears in human mode}\label{sec:trouble-human}

\paragraph{Symptom.} With \code{render\_mode="human"} nothing is shown, and
a \code{RuntimeWarning} states that the active matplotlib backend is
non-interactive.

\paragraph{Cause.} \code{"human"} mode opens a window through matplotlib,
which needs an interactive backend such as TkAgg or QtAgg and a display.
When matplotlib falls back to Agg (because no GUI toolkit is installed, no
display is available, or \code{MPLBACKEND=Agg} is set), the environment warns
once and renders nothing. Pistonball's window is opened by pygame and needs a
display, but no matplotlib backend.

\paragraph{Solution.} Check which backend is active:

\begin{shellcode}
python -c "import matplotlib; print(matplotlib.get_backend())"
\end{shellcode}

If it prints \code{agg} (any case), install a GUI toolkit and select it:

\begin{itemize}
  \item TkAgg needs Tk: \code{sudo apt-get install python3-tk} (Debian,
    Ubuntu), \code{sudo dnf install python3-tkinter} (Fedora),
    \code{brew install python-tk} (Homebrew) or \code{conda install tk}.
  \item QtAgg needs a Qt binding: \code{pip install PyQt6} (or \code{PySide6}).
\end{itemize}

Then unset \code{MPLBACKEND} or set it explicitly, for example
\code{export MPLBACKEND=TkAgg}. In a Jupyter notebook, \code{"human"} mode is
not the right tool: render with \code{"rgb\_array"} and display the frames
with \code{matplotlib.pyplot.imshow} or save an animation. Over SSH, either
enable X forwarding (\code{ssh -X}) or record the episode and copy the file.

\section{Missing optional dependencies}\label{sec:trouble-deps}

Optional dependencies are imported only when a feature needs them, and the
resulting error names the extra that provides them
(Table~\ref{tab:trouble-deps}).

\begin{table}[htbp]
  \centering
  \small
  \caption{Errors caused by missing optional dependencies.}
  \label{tab:trouble-deps}
  \begin{tabularx}{\textwidth}{@{}>{\raggedright\arraybackslash}X>{\raggedright\arraybackslash}p{4.4cm}@{}}
    \toprule
    Error message (abbreviated) & Fix \\
    \midrule
    \code{PistonballEnv requires pymunk (and pygame for rendering)} &
      \code{pip install ".[pistonball]"} \\
    \code{Pistonball rendering requires pygame} & \code{pip install ".[pistonball]"} \\
    \code{toolkit.neural\_toolkit requires PyTorch} & \code{pip install ".[torch]"} \\
    \code{No module named 'torch'} when creating \code{KuramotoOscillatorEnvTorch}
      or a \envid{KuramotoOscillatorTorch-*} id & \code{pip install ".[torch]"} \\
    \code{Video export requires imageio and imageio-ffmpeg} & \code{pip install ".[video]"} \\
    \code{The parakit GUI requires Tk (the 'tkinter' module)} &
      system package \code{python3-tk} or \code{conda install tk}
      (Section~\ref{sec:install-tk}) \\
    \code{No module named 'pytest'} or \code{'ruff'} & \code{pip install -e ".[all,dev]"} \\
    \bottomrule
  \end{tabularx}
\end{table}

Note that \code{pip install ".[dev]"} installs only the development tools; the
test suite skips the tests of optional features that are not installed, so a
complete test run needs \code{".[all,dev]"}. If an extra is installed but the
error persists, the package was probably installed into a different Python
environment than the one running the code; compare \code{which python} with
\code{pip -V}, and use \code{python -m pip install ...} to be sure.

\section{``Out of date'' warnings for -v0 environment ids}\label{sec:trouble-outofdate}

\paragraph{Symptom.} \code{gymnasium.make("WirelessComm-v0")} or
\code{gymnasium.make("KuramotoOscillator-v0")} prints

\begin{textcode}
DeprecationWarning: WARN: The environment WirelessComm-v0 is out of date.
You should consider upgrading to version `v1`.
\end{textcode}

\paragraph{Cause and solution.} Gymnasium assumes that a higher version
suffix is a newer revision of the same environment. In \pkg{env\_lib} the
suffixes denote preset configurations: \envid{WirelessComm-v1} is the
$4 \times 4$ grid, not an improved \envid{WirelessComm-v0}. The warning is
harmless and can be ignored. \code{env\_lib.make}, which otherwise behaves
like \code{gymnasium.make}, suppresses it. To silence it when calling
\code{gymnasium.make} directly:

\begin{pycode}
import warnings
import gymnasium as gym
import env_lib  # registers the ids

warnings.filterwarnings("ignore", message=r".*is out of date", category=DeprecationWarning)
env = gym.make("WirelessComm-v0")
\end{pycode}

\section{AJLATT deprecation warnings}\label{sec:trouble-ajlatt}

\paragraph{Symptom.} Creating an AJLATT environment prints warnings such as
\code{AJLATT parameter 'num\_Robot' is deprecated; use 'num\_robots'} or
\code{parameter 'figID' no longer has any effect and is ignored}.

\paragraph{Cause.} The environment is configured by the \code{AJLATTConfig}
dataclass since version 1.0. The legacy keyword names of the former
\pkg{argparse}-based configuration are still accepted by
\code{AJLATTEnv(...)}, \code{env\_lib.ajlatt\_env(...)},
\code{env\_lib.ajlatt\_env.make(...)} and \code{AJLATTConfig.from\_kwargs},
but translated with a \code{DeprecationWarning}.

\paragraph{Solution.} Rename the arguments as in
Table~\ref{tab:trouble-ajlatt}; Appendix~\ref{app:migration} shows complete
before-and-after examples.

\begin{table}[htbp]
  \centering
  \small
  \caption{Legacy AJLATT parameters and their replacements.}
  \label{tab:trouble-ajlatt}
  \begin{tabularx}{\textwidth}{@{}>{\raggedright\arraybackslash}p{3.8cm}>{\raggedright\arraybackslash}X@{}}
    \toprule
    Legacy argument & Replacement \\
    \midrule
    \code{num\_Robot}, \code{num\_Target} & \code{num\_robots}, \code{num\_targets} \\
    \code{T\_steps} & \code{max\_episode\_steps} \\
    \code{maxmum\_run} & \code{max\_episode\_steps}; the episode is limited to \code{maxmum\_run + 1} steps \\
    \code{SIGPERCENT} & \code{range\_noise\_proportional} \\
    \code{useupdate} & \code{use\_update} \\
    \code{fix\_seed} & \code{seed=<int>} in the configuration, or \code{env.reset(seed=...)} \\
    \code{render=True} & \code{render\_mode="human"} (or \code{"rgb\_array"}) \\
    \code{figID}, \code{record}, \code{ros}, \code{log\_dir}, \code{episode\_num}, \code{env\_name}, \ldots &
      no effect; remove them \\
    \bottomrule
  \end{tabularx}
\end{table}

Unknown names raise a \code{TypeError} in \code{AJLATTEnv(...)} and
\code{AJLATTConfig.from\_kwargs}, and are ignored with a warning by the
backward-compatible factory \code{make}. Python hides
\code{DeprecationWarning} outside the \code{\_\_main\_\_} module by default;
run \code{python -W default::DeprecationWarning script.py} to see every
translation when migrating a code base.

\section{PyTorch, GPUs and devices}\label{sec:trouble-torch}

\paragraph{Installing a CPU or a CUDA build.}
The \code{torch} extra installs the default PyTorch wheel of the platform,
which on Linux is a large CUDA build. For CPU-only machines install the CPU
wheel first (\code{pip install torch --index-url
https://download.pytorch.org/whl/cpu}); for a specific CUDA version use the
command shown at \url{https://pytorch.org/get-started/}. Check the result
with

\begin{shellcode}
python -c "import torch; print(torch.__version__, torch.cuda.is_available())"
\end{shellcode}

\paragraph{The Kuramoto PyTorch backend.}
\code{KuramotoOscillatorEnvTorch} accepts \code{device="cpu"},
\code{"cuda"}, \code{"cuda:<index>"}, a \code{torch.device}, or
\code{"auto"}, which selects the first GPU when CUDA is available and the CPU
otherwise (the registered id \envid{KuramotoOscillatorTorch-v2} uses
\code{"auto"}). Requesting \code{"cuda"} on a machine without CUDA raises a
\code{ValueError} that suggests \code{"auto"} or \code{"cpu"}. Observations
returned by \code{reset} and \code{step} are always NumPy arrays; the batched
accessors \code{get\_batch\_observations()} and \code{get\_batch\_rewards()}
return tensors on the environment's device. A GPU pays off for large batches
of systems; for a single small system the CPU is faster because of the
transfer and launch overhead.

\paragraph{neural\_toolkit devices.}
The networks take a \code{device} argument (default \code{"cpu"}) that is
passed to \code{Module.to}; \code{"auto"} is not accepted there. Choose the
device explicitly and move the inputs to it:

\begin{pycode}
import torch
from toolkit.neural_toolkit import MLPPolicyNetwork

device = "cuda" if torch.cuda.is_available() else "cpu"
policy = MLPPolicyNetwork(8, 4, device=device)
logits = policy(torch.randn(32, 8, device=device))
\end{pycode}

\code{NetworkUtils.load\_checkpoint} loads tensors onto the device of the
model by default, so checkpoints written on a GPU load on a CPU-only machine.
Checkpoints written before version 1.0 cannot be loaded
(Section~\ref{sec:neural-changes}).

\section{Video export}\label{sec:trouble-video}

\paragraph{Symptom.} Saving a \code{.mp4} (or \code{.webm}, \code{.avi},
\code{.mov}, \code{.mkv}) file fails with \code{Video export requires imageio
and imageio-ffmpeg}, or \pkg{imageio} reports that it cannot find an
\code{ffmpeg} executable.

\paragraph{Solution.} Install the \code{video} extra,
\code{pip install ".[video]"}. The \pkg{imageio-ffmpeg} package ships its own
\code{ffmpeg} binary, so no system installation is needed. GIF export uses
Pillow, which is always installed with matplotlib, and is the fallback when a
video codec is unavailable. \code{save\_animation} rejects other suffixes with
a \code{ValueError}; frames must all have the same shape
\code{(H, W, 3)} or \code{(H, W, 4)}.

\section{Reproducibility and seeding}\label{sec:trouble-seeding}

The environments are deterministic given their seeds and never use global
random state, but a complete experiment contains several independent sources
of randomness. Seed each of them:

\begin{itemize}
  \item \emph{Environment.} Call \code{env.reset(seed=s)} once at the start of
    a run; later calls to \code{env.reset()} without a seed continue the same
    random stream, so a whole sequence of episodes is reproducible. Resetting
    with the same seed every episode repeats the same initial state.
  \item \emph{Action sampling.} \code{env.action\_space.sample()} uses the
    space's own generator, which \code{reset} does not seed. Call
    \code{env.action\_space.seed(s)}. This also applies to
    \code{record\_episode} without a policy.
  \item \emph{Construction-time randomness.} Some environments draw random
    structure in the constructor, independent of \code{reset}: the
    Erd\H{o}s-R\'enyi graph of \code{ConsensusEnv} (\code{graph\_seed}; with
    the default \code{None} two instances may differ, even in their
    observation size) and the random topology of the Kuramoto environments
    (\code{topology\_seed=42}). AJLATT accepts a default seed in its
    configuration (\code{seed=...}), used by the first \code{reset()} without
    a seed.
  \item \emph{Networks.} Call \code{torch.manual\_seed(s)} before creating
    networks and before training. \code{VariationalEncoder} samples its noise
    from PyTorch's default generator. On GPUs some operations are
    non-deterministic; see the PyTorch documentation on reproducibility
    (\code{torch.use\_deterministic\_algorithms}).
  \item \emph{Tabular tools.} Pass \code{rng=s} to the tables; each table
    owns its generator.
  \item \emph{Your own code.} Use a \code{numpy.random.default\_rng(s)}
    generator instead of the global \code{np.random} functions.
\end{itemize}

Seeded trajectories are reproducible for a fixed version of the package and
its dependencies. Across versions, results can change where bugs were fixed;
Appendix~\ref{app:migration} lists these changes. The AJLATT environment can
reproduce its behaviour before version 1.0 closely with
\code{ci\_solver="slsqp"}.

\section{Performance}\label{sec:trouble-performance}

\begin{itemize}
  \item \emph{Render only when needed.} Rendering a frame costs 3 to 36\,ms,
    which is far more than a simulation step (Table~\ref{tab:examples-bench}).
    Train with \code{render\_mode=None}, and record evaluation episodes
    separately, with \code{render\_every} greater than one for long episodes.
    \code{"human"} mode additionally limits the frame rate, so never train in
    \code{"human"} mode.
  \item \emph{Create environments once.} Construction loads maps, builds
    graphs and allocates buffers; reuse an environment across episodes with
    \code{reset()} instead of creating a new one per episode. The first
    \code{render()} call builds the figure and is slower than later calls.
  \item \emph{Choose the integrator.} For the Kuramoto environments, Euler
    integration is about twice as fast as RK4; use RK4 when accuracy at large
    \code{dt} matters.
  \item \emph{Batch on the GPU.} Many Kuramoto systems simulated together by
    \code{KuramotoOscillatorEnvTorch} (\code{n\_agents} systems per step) are
    much cheaper than separate environments, especially on a GPU.
  \item \emph{AJLATT.} Keep the default \code{ci\_solver="newton"}, which is
    several times faster than \code{"slsqp"}; the cost per step grows with the
    number of robots because every robot fuses its neighbours' estimates.
  \item \emph{Small networks on the CPU.} PyTorch uses several threads by
    default. For the small networks typical of these environments, and when
    several training processes share a machine, thread contention can make
    training much slower; limit the threads with \code{torch.set\_num\_threads(1)}
    or the environment variable \code{OMP\_NUM\_THREADS=1} and compare.
  \item \emph{Measure.} \code{benchmarks/benchmark\_envs.py}
    (Section~\ref{sec:examples-benchmarks}) measures step and frame times on
    your machine; \code{python -m cProfile -s cumtime script.py} shows where
    time goes in your own code.
\end{itemize}

\section{Other common errors}\label{sec:trouble-other}

\begin{description}[leftmargin=0.6cm,style=nextline]
  \item[\code{step()} raises before the first \code{reset()}]
    Every environment requires \code{reset()} before the first \code{step()}
    and raises \code{RuntimeError} (\code{gymnasium.error.ResetNeeded} for the
    Kuramoto environments) otherwise. Environments created with
    \code{gymnasium.make} raise Gymnasium's own error from the
    \code{OrderEnforcing} wrapper.
  \item[An episode is longer or shorter than expected]
    Each environment enforces its own episode limit through a constructor
    argument (\code{max\_steps}, \code{max\_iter}, \code{max\_cycles} or
    \code{max\_episode\_steps}); pass it to the constructor or to
    \code{env\_lib.make}. Passing \code{max\_episode\_steps} to
    \code{gymnasium.make} adds a second \code{TimeLimit} wrapper, and the
    shorter of the two limits applies.
  \item[\code{AttributeError} for an environment method]
    \code{gymnasium.make} and \code{env\_lib.make} return a wrapped
    environment. Access environment-specific methods and attributes through
    \code{env.unwrapped}, for example
    \code{env.unwrapped.laplacian\_policy()}.
  \item[\code{render()} returns \code{None} with a warning]
    The environment was created without a render mode. Pass
    \code{render\_mode="rgb\_array"} (or \code{"human"}) to the constructor or
    to \code{make}.
  \item[Actions are rejected]
    Actions must match \code{env.action\_space} in shape. Pistonball in
    discrete mode expects one value in $\{0, 1, 2\}$ per piston, not a single
    integer; AJLATT expects an array of shape \code{(num\_robots, 2)}.
    Appendix~\ref{app:migration} describes the changes from earlier versions.
  \item[The \pkg{parakit} window does not open]
    \code{ParameterTuner.tune()} raises \code{ImportError} without Tk and
    \code{RuntimeError} without a display; use the headless API
    (Section~\ref{sec:parakit-gui}).
  \item[Figures look different from the examples]
    \pkg{plotkit} applies its style locally and does not change global
    settings. If other code changed \code{matplotlib.rcParams} (or called
    \code{set\_research\_style} without \code{reset\_style}), the preset is
    applied on top of those settings; call \code{matplotlib.rcdefaults()} to
    start from the matplotlib defaults.
\end{description}

\par\endgroup
